\documentclass[a4paper,12pt]{article}

% --- Use XeLaTeX font support ---


\usepackage{fontspec}
\setmainfont{TeX Gyre Termes} % nice Times-like font, installed by texlive-core
\setsansfont{TeX Gyre Heros} % optional sans-serif
\setmonofont{Fira Code}       % optional monospaced font
% --- Math packages ---
\usepackage{amsmath, amssymb, amsthm}
\usepackage{mathtools}

% --- Page layout ---
\usepackage[margin=1in]{geometry}

% --- Better lists ---
\usepackage{enumitem}
% --- Hyperlinks for references ---
\usepackage[colorlinks=true, linkcolor=blue, urlcolor=blue]{hyperref}

% Header: fancy layout like the rovided image
\usepackage{fancyhdr}
\setlength{\headheight}{26pt} % increase if fancyhdr warns
\pagestyle{fancy}
\fancyhf{} % clear header/footer
\rhead{\begin{tabular}{r} Mt.: 3052737\end{tabular}}
\chead{\begin{tabular}{c}Lösungen EA 2\end{tabular}}
\lhead{\begin{tabular}{l}61211 Analysis SS26 \end{tabular}}
\renewcommand{\headrulewidth}{0.8pt}
\fancyfoot[C]{- \thepage\ -}
% Ensure the plain page style (used by \maketitle) uses the same header

% --- Line spacing ---
\usepackage{setspace} % provides \doublespacing and \onehalfspacing

% --- Optional solutions environment ---
\newenvironment{solution}{
  \par\noindent\textbf{Solution:}\quad
}{\hfill$\blacksquare$\par\vspace{1em}}

% % --- Title info ---
% \title{Aufgabe \#1}
% \author{Your Name}
% \date{\today}
% 
\begin{document}
\raggedright
% Enable double (2.0) line spacing for the document body
\doublespacing


\noindent\textbf{Aufgabe 2} \\
\noindent\textbf{(a)} Die Aussage ist wahr. 
Die Funktion springt bei $-\sqrt 2$ von 1 auf 0 und bei $\sqrt 2$ von 0 auf 1. Da  $\sqrt 2$ nicht zum Definitionsbereich gehört kann sie an diesen Stellen keinen Unstetigkeitspunkt haben. \\
Also
\begin{itemize}
    \item entweder rationales $q \in (- \sqrt 2, \sqrt 2)$ mit $f(q) =0$, dann gibt es mit genügend kleinem $\epsilon > 0$ ($ \epsilon < |q - \sqrt 2|$) eine Umgebung $U_{\epsilon}(q)$, so dass alle rationalen $x \in U_{\epsilon}(q)$ ebenfalls in $ (- \sqrt 2, \sqrt 2)$ liegen mit $f(x) = 0$
    \item oder, vollkommen analog, rationales $q \notin (- \sqrt 2, \sqrt 2)$ dann lässt sich eine Umgebung finden, sodass $f(q) = f(x) =1$ ist.
\end{itemize}

\noindent\textbf{(b)}
Die Funktionen $g,h$ sind stetig. \\
 Intuitiv: Ein eindimensionaler Schnitt durch eine stetige Fläche kann keine 'problematischen' Stellen enthalten. \\
Formal:
Man definiert eine Hilfsfunktion $\varphi: \mathbb{R} \to \mathbb{R}^2$ mit $$\varphi(y) = \begin{pmatrix} a \\ y \end{pmatrix}$$ \\
Die Funktion $\varphi$ ist offensichtlich stetig. \\
Nun ist $g(y) = f(a,y)= f(\varphi(y))$ und es greift die Kettenregel für stetige Funktionen, d.h. $g = f \circ \varphi $ ist ebenfalls stetig. \\
Für die Funktion $h$ vollkommen analog.



\end{document}